Modern multifunction phased-array radars must schedule heterogeneous missions while facing coordinated, adaptive interference. Link-level signal-to-interference ratios do not by themselves reveal whether a radar network completes its deadline-constrained mission load. We introduce FluxPhased, a physics-grounded task-level benchmark for adversarial co-adaptation between multifunction radars and cooperative jammers. Missions carry service and bearing identities, deadlines, and exact detection credit; the link model combines generalized planar-array responses with finite jammer energy budgets. A pre-training contestability test and a four-view evaluation protocol separate natural scheduling loss from adversarial loss.

In a one-jammer versus two-radar game, the valid 12-dB S6 runs neutralize $63.7\%\pm1.0\%$ of the jammer's floor-adjusted disruptive power across two training seeds. When the attacker team is doubled under the same 63-token team budget, three converged seeds yield a stable 2v2 equilibrium with head-to-head mission drop $0.3366\pm0.0143$ and neutralization $23.0\%\pm1.1\%$. Radar competence against idle jammers remains low, showing that the degradation is an offensive scaling effect rather than a corresponding defender collapse. A co-located-jammer control attributes most of the change to attacker count under the S7 separated-radar geometry (neutralization $28.4\%$), with cross-fire geometry contributing a further reduction to $24.2\%$. Finally, a lightweight opponent-class mixture reduces singleton relative leverage from $0.502$ to $0.262$ at 75\% singleton exposure, with the largest neutralization value at 50\% exposure. The results establish a measurable separation between equilibrium stability, attack scaling, and opponent-distribution robustness.
